%%%%%%%%%%%%%%%%%%%%%%%%%%%%%%%%%%%%%%%%%%%%%%%%%%%%%
%			HLAVIČKA								%
%%%%%%%%%%%%%%%%%%%%%%%%%%%%%%%%%%%%%%%%%%%%%%%%%%%%%
\documentclass[openany,12pt]{memoir}

\usepackage[utf8]{inputenc} 
\usepackage[czech]{babel}
\usepackage[T1]{fontenc}
\usepackage[top=1.5cm, bottom=2cm, left=2cm, right=2cm]{geometry}  % --> NASTAVENÍ OKRAJŮ
\usepackage{fancyhdr}
\usepackage{graphicx}
%\usepackage{lmodern}
\usepackage{tgschola}  % --- Silnější font
\usepackage{xwatermark}
\usepackage{xcolor}
\usepackage{changepage}
\usepackage{newunicodechar}



%%%%%% Package na zpěvník
\usepackage[full]{leadsheets}%http://mirrors.nic.cz/tex-archive/macros/latex/contrib/leadsheets/leadsheets_en.pdf   --> dokumentace	
\definesongtitletemplate{empty}{} 
\setchords{
format = \bfseries,   %tučné akordy
minor = {mi},% 
input-notation = {german},%
output-notation = {german}%
}
\definesongtitletemplate{empty}{} 

\newlength{\drop}
\newwatermark[allpages,color=red!50,angle=0,scale=2, xpos=0,ypos=0]{\includegraphics[width=5cm]{obr/pozadi2.jpg}} %--> dvojka na pozadí


%%%% Vlastní příkazy
\newcounter{Slokočet}   %Automatické číslování slok
\newcommand{\mezera}{\vspace*{0.5cm}}   %Horizontální odsazení slok
\newcommand{\stred}{5.2cm}   %%% Na zarovnání slok doprostřed, pozn. automatičtější zarovnávání na střed nejde
\newcommand{\refren}{\mezera \noindent \textbf{R:} } %refrén
\newcommand{\sloka}{\mezera \noindent \addtocounter{Slokočet}{1} \arabic{Slokočet}. } 	%sloka, která se automaticky čísluje
\newcommand{\ssloka}{\mezera \noindent}  % vlastní číslo sloky
\newcommand{\carka}{,\:}
\newcommand{\predehra}{\mezera \hspace*{-1.85cm} \textbf{Předehra:} }
\newcommand{\m}[1]{\color{white}{#1}}  %Pro akordy
\newcommand{\ap}{'} %Pro apostrof
\newcommand{\elipsa}{\kern\fontdimen3\font} %Příkaz pro lepší zacházení s výpustkami (=...); je to vpodstatě jen mezera mezi tečkama výpustky

%%%%%%%%%%%%%%%%%%%%%%%%%%%%%%%%%%%%%%%%%%%%%%%%%%%%%%
%			NÁVOD									 %
%%%%%%%%%%%%%%%%%%%%%%%%%%%%%%%%%%%%%%%%%%%%%%%%%%%%%%
%1. Věci v hlavičce IGNOROVAT
%2. Píseň psát do prostoru mezi \begin{song} a \end{song}
%3. další řádek se značí dvěma odsazeníma (= dvakrát stisknout enter)
%4. \refren vždy na začátku refrenu a \sloka na začátek sloky (automaticky se čísluje)
%  \ = alt gr + q ; [/] = alt gr f/g ; {/} = alt gr + b/n; ^ = alt gr + 3 + mezera
%Cokoliv napíšete do ^{  } se bude brát jako akord
%když se toto bude dotýkat nějakého slova (nebude mezi tím a slovem mezera)
%tak se akordy zjeví nad slovem, ale pište to před slovo
%Když se to nedotýká slova, tak akord lítá ve vzduchu a vytiskne se větší mezera
%První možnost je asi preferovanější
%5. Akordy stačí psát jen do první sloky, když se nezmění -- kytaristi to zvládnou



%\usepackage{subfiles}
%</++>

\begin{document}

\pagestyle{empty}
%\addtocontents{toc}{\protect\thispagestyle{empty}}
%\tableofcontents \thispagestyle{empty}\newpage
\newgeometry{top=1.5cm, bottom = 0cm, left = 2cm, right = 2cm}
	
\newpage

%Pro kompilaci je důležitý, aby ten \input a \end{document} zůstali na posledních
%dvou řádcích!!!!


\begin{song}{title=\centering Slaboch Ben \\\normalsize Kapitán Kid  \vspace*{-0.3cm}}  %% sem se napíše jméno songu a autor
\moveright 0.8cm \vbox{      %Varianta č. 1  ---> Jeden sloupec zarovnaný na střed	
\begin{minipage}[t]{0.48\textwidth}\setlength{\parindent}{0.45cm}  %Varianta č. 2 --> Dva sloupce

\sloka 
	^{C}V lesích ^{Ami}Ontaria ^{F}čerstvá smůla ^{G}voní.

	^{C}Tam, kde silných ^{A7}mužů paže ^{Dmi}vládne ^{G}jen,

	kde ^{C}sosny ^{Ami}staleté se ^{F}pod sekyrou ^{Cdim}kloní,

	v kempu ^{C}dřevorubců ^{G7}žije Slaboch ^{C}Ben.
	
\sloka
^{F}Mlčí stále ^{Cdim}jen a ^{C}snáší strasti ^{Ami}žití

a když ^{Dmi}kamarádů ^{G7}výsměch dávno ^{Ami}ztich.

^{E}V noci nad ^{Ami}jezerem, ^{Dmi}Polárka když ^{Ami}svítí,

^{E}odpouští jim věčný jejich smích.

\refren
	^{A}Že je slaboch, ^{A7}že prý nemá

	^{D}sílu, že se vrány bojí,

	^{E}říká o něm ^{E7}předák Andy ^{A}Renn.

	(Ten kojot) Hodí se jen ^{A7}k tomu, aby

	^{D}chodil kempu ^{D7}pro zásoby,

	^{E}je prý jenom ^{E7}pro smích Slaboch ^{A}Ben. ^{G7}

\sloka
	Na cestu mu dává bez nábojů zbraně,
	
	kdyby se snad před kojoty bránit chtěl.
	
	Kdyby se mu chtělo vystřelit si na ně,
	
	že by stejně, jak se míří nevěděl.
	
	Bez nábojů prý se aspoň nepostřelí,
	
	kdyby třeba přemoh strach a ránu dal.
	
	A když ještě dodal: "Snad se vrátíš celý,"
	
	kemp se málem smíchy rozsypal.
	
\refren
\end{minipage}\begin{minipage}[t]{0.48\textwidth}\setlength{\parindent}{0.45cm}\vspace*{0.59cm}  % V případě varianty č.2 jde odsud text do pravé části

\sloka
	Byli by se smáli ještě hodnou chvíli,
	
	kdyby tábor nebyl náhle přepaden
	
	bandou, kterou vedl jednooký Billy,
	
	kterým mnohý traper byl už zastřelen.
	
	Na pět kamarádů deset koltů cílí,
	
	každý z nich už s životem se rozžehnal.

	Když tu náhle výkřik zazní v pravou chvíli:
	
	\uv{Kdo se nevzdá, žít nebude dál!}
	
\refren
	To že je slaboch? Ten že nemá
	
	sílu, že se rány bojí,
	
	říká si teď v duchu Andy Renn.
	
	(Ten kojot) Ten že se nám hodí k tomu,
	
	aby chodil pro zásoby,
	
	tohle přece není slaboch Ben.

\sloka
	Deset koltů nato cíl svůj ihned mění,
	
	míří tam, kde v jejich zádech stojí Ben.

	Leč pět kamarádů v tomtéž okamžení
	
	střelbu spustí, dokud Bill je překvapen.
	
	Pak už v trávě leží jednooký Billy
	
	a s ním jeho čtyři muži mrtvi jen.
	
	Jako šestý s nimi s nenabitou zbraní,
	
	pro své pardy zemřel Slaboch Ben.

\refren
	Že je slaboch, že prý nemá
	
	sílu, že se vrány bojí,
	
	říkal o něm kdysi Andy Renn.
	
	Na hrobě, jenž teď ho kryje
	
	nápis je, že proto žije
	
	pardů pět, že zemřel pro ně Ben.
\end{minipage}
}
\setcounter{Slokočet}{0}
\end{song}


\end{document}	